\documentclass[UTF8,zihao=-4,a4paper,oneside]{ctexbook}

% ============================================================
% 基础宏包
% ============================================================

\usepackage{amsmath}
\usepackage{amssymb}
\usepackage{amsthm}
\usepackage{mathtools}

\usepackage{geometry}
\usepackage{graphicx}
\usepackage{tikz}
\usepackage{xcolor}
\usepackage{fontspec}

\usepackage{titlesec}
\usepackage{titletoc}

\usepackage{fancyhdr}
\usepackage{lastpage}

\usepackage{enumitem}
\usepackage{booktabs}
\usepackage{array}
\usepackage{multirow}

\usepackage{minted}

\usepackage{hyperref}
\usepackage{bookmark}

% ============================================================
% 页面设置
% ============================================================

\geometry{
    a4paper,
    left=24mm,
    right=24mm,
    top=24mm,
    bottom=24mm,
    headheight=16pt
}

% ============================================================
% 颜色主题
% ============================================================

\definecolor{MainColor}{RGB}{38, 70, 83}
\definecolor{AccentColor}{RGB}{42, 157, 143}
\definecolor{LightColor}{RGB}{242, 247, 247}

% VSCode Light Modern 配色
\definecolor{VsBg}{HTML}{f7f7f0}       % 编辑器背景
\definecolor{VsText}{HTML}{3B3B3B}      % 普通文本
\definecolor{VsKeyword}{HTML}{0000FF}   % 类型/声明关键字（蓝）
\definecolor{VsControl}{HTML}{AF00DB}   % 控制流/预处理（紫）
\definecolor{VsComment}{HTML}{008000}   % 注释（绿）
\definecolor{VsString}{HTML}{A31515}    % 字符串（红）
\definecolor{VsNumber}{HTML}{098658}    % 数字（青绿）
\definecolor{VsFunction}{HTML}{795E26}  % 函数名（黄褐）
\definecolor{VsClass}{HTML}{267F99}     % 类名（青）
\definecolor{VsLineNo}{HTML}{6E7681}    % 行号（灰）
\definecolor{VsFrame}{HTML}{CFCFCF}     % 边框（浅灰）

% ============================================================
% 超链接
% ============================================================

\hypersetup{
    colorlinks=true,
    linkcolor=MainColor,
    citecolor=AccentColor,
    urlcolor=AccentColor,
    bookmarksopen=true,
    bookmarksnumbered=true,
    pdftitle={XCPC 算法竞赛模板},
    pdfauthor={gujiu6}
}

% ============================================================
% 代码环境（VSCode Light Modern 风格 + Fira Code 字体）
% 基于 minted（Pygments），编译需加 -shell-escape：
%   xelatex -shell-escape main.tex
% 配色由 _minted/vslight.style.minted 定义（勿删除，否则回退默认配色）。
% ============================================================

% Fira Code 安装在用户字体目录，按文件名加载
\IfFontExistsTF{Fira Code}{\setmonofont{Fira Code}}{\setmonofont{Latin Modern Mono}}

\usemintedstyle{vslight}

\setminted{
    fontsize=\small,
    breaklines=true,
    breakanywhere=true,
    linenos=true,
    numbersep=10pt,
    frame=single,
    framerule=0.4pt,
    rulecolor=VsFrame,
    bgcolor=VsBg,
    tabsize=4
}

% 生成 cpp 环境，样式与 listings 对齐（白底、浅灰边框、行号）
\newminted[cpp]{cpp}{
    fontsize=\normalsize,
    breaklines=true,
    breakanywhere=true, % minted 把空格包成 \PYG{+w} 宏，无空格的长 token 串（如 G2[u].erase(std::unique(...))）
                        % 找不到断点而溢出，breakanywhere 允许在任意字符处换行兜底
    linenos=true,
    numbersep=10pt,
    frame=single,
    framerule=0.4pt,
    rulecolor=VsFrame,
    bgcolor=VsBg,
    tabsize=4,
    % 把自定义类型名加入 extrakeywordsdeclaration（渲染成 VsClass 青，与
    % struct X 声明处的 Name.Class 一致）。Pygments 无法做语义类型追踪，
    % 只能枚举：新增模板类型时，把类型名加进下面这个大括号即可。
    % 注意：只能写单个标识符（如 vector），不能写 std::vector——
    % Pygments 会把 std、::、vector 拆成多个 token，写在一起永远匹配不上。
    extrakeywordsdeclaration={
        Comp,Complex,FWT,MInt,ModInt,Poly,Z,arbModInt,
        i128,i64,p32,p64,u32,u64,
        Disjoint,DisjointWithUndo,FHQTreap,MergeSplitSegmentTree,MinCostFlowHLPP,PersistantFHQTreap,
        SparseTable,Splay,_KV
    }
}

% 让所有章节中已有的 \begin{lstlisting} 直接复用 minted
\newenvironment{lstlisting}
  {\VerbatimEnvironment\begin{cpp}}
  {\end{cpp}}

% 行号默认被 fancyvrb 硬编码成 \rmfamily\tiny（不随代码字号缩放）。
% 这里改为跟随代码块字体（fontsize 设为多少行号就多大），并用 VsLineNo 灰。
\renewcommand{\theFancyVerbLine}{\textcolor{VsLineNo}{\arabic{FancyVerbLine}}}

% ============================================================
% 数学环境
% ============================================================

\newtheoremstyle{algstyle}
    {8pt}
    {8pt}
    {\normalfont}
    {}
    {\bfseries\color{MainColor}}
    {.}
    {0.5em}
    {}

\theoremstyle{algstyle}

\newtheorem{theorem}{定理}[chapter]
\newtheorem{lemma}[theorem]{引理}
\newtheorem{proposition}[theorem]{命题}
\newtheorem{corollary}[theorem]{推论}

\theoremstyle{definition}

\newtheorem{definition}[theorem]{定义}
\newtheorem{example}[theorem]{例}

% ============================================================
% 标题样式
% ============================================================

\titleformat
{\chapter}
[display]
{\normalfont\bfseries\color{MainColor}}
{\filleft\fontsize{18}{18}\selectfont 第\,\thechapter\,章}
{5pt}
{\filright\fontsize{32}{38}\selectfont}
[\vspace{5pt}{\titlerule[1.2pt]}]

\titleformat
{\section}
{\Large\bfseries\color{MainColor}}
{\thesection}
{0.8em}
{}

\titleformat
{\subsection}
{\large\bfseries\color{AccentColor}}
{\thesubsection}
{0.8em}
{}

% ============================================================
% 目录样式
% ============================================================

% 章标签用中文数字（第十一章为 4 字），盒子宽度需容纳 4 个汉字，
% 否则“第十一章”溢出盒子与标题重叠（chapter 超过 10 后出现）。
\titlecontents{chapter}
[3.25em]
{\bfseries}
{\contentslabel{4.6em}}
{}
{\titlerule*[0.7pc]{.}\contentspage}

% 节编号两位（如 11.10）时盒子需加宽，否则“11.10”溢出盒子。
\titlecontents{section}
[4.75em]
{}
{\contentslabel{2.4em}}
{}
{\titlerule*[0.7pc]{.}\contentspage}

\titlecontents{subsection}
[6em]
{}
{\contentslabel{2.5em}}
{}
{\titlerule*[0.7pc]{.}\contentspage}

\titlecontents{subsubsection}
[7.25em]
{}
{\contentslabel{3em}}
{}
{\titlerule*[0.7pc]{.}\contentspage}

% 目录页码框默认 1.55em，粗体 3 位页码（如第十一章的 108）会略溢出；
% 加宽到 1.8em（页码右对齐，视觉位置不变）。注：内核中 \@pnumwidth 是宏，
% 不是长度寄存器，需用 \def 重定义而非 \setlength。
\makeatletter
\def\@pnumwidth{1.8em}
\makeatother

% ============================================================
% 页眉页脚
% ============================================================

\pagestyle{fancy}

\fancyhf{}

\fancyhead[L]{
    \small
    \textcolor{MainColor}{\textbf{算法竞赛模板}}
}

\fancyhead[R]{
    \small
    \textcolor{gray}{\nouppercase{\leftmark}}
}

\fancyfoot[C]{
    \small
    \textcolor{gray}{\thepage\ /\ \pageref{LastPage}}
}

\renewcommand{\headrulewidth}{0.5pt}
\renewcommand{\footrulewidth}{0pt}

% 章节首页（plain 样式）也显示页脚页码
\fancypagestyle{plain}{
    \fancyhf{}
    \fancyfoot[C]{
        \small
        \textcolor{gray}{\thepage\ /\ \pageref{LastPage}}
    }
    \renewcommand{\headrulewidth}{0pt}
    \renewcommand{\footrulewidth}{0pt}
}

% ============================================================
% 封面 Logo
% ============================================================

\newcommand{\AlgorithmLogo}{
    \begin{tikzpicture}[scale=1]
        % 外圈
        \draw[
            MainColor,
            line width=2pt
        ]
        (0,0) circle (1.3);

        % 节点
        \fill[AccentColor] (0,0.75) circle (0.12);
        \fill[AccentColor] (-0.75,-0.4) circle (0.12);
        \fill[AccentColor] (0.75,-0.4) circle (0.12);

        % 连线
        \draw[
            MainColor,
            line width=1.5pt
        ]
        (0,0.75) -- (-0.75,-0.4);

        \draw[
            MainColor,
            line width=1.5pt
        ]
        (0,0.75) -- (0.75,-0.4);

        \draw[
            MainColor,
            line width=1.5pt
        ]
        (-0.75,-0.4) -- (0.75,-0.4);

        % 中心节点
        \fill[MainColor] (0,0) circle (0.18);
    \end{tikzpicture}
}

% ============================================================
% 封面
% ============================================================

\newcommand{\MakeAlgorithmTitle}{
    \begin{titlepage}

        \centering

        \vspace*{2cm}

        % Logo
        \IfFileExists{logo.png}{
            \includegraphics[width=4cm]{logo.png}
        }{
            \AlgorithmLogo
        }

        \vspace{1.5cm}

        {\fontsize{38}{45}\selectfont
        \bfseries
        \textcolor{MainColor}{XCPC Templates}
        }

        \vspace{0.5cm}

        {\fontsize{20}{28}\selectfont
        \textcolor{AccentColor}{常用 XCPC 模板}

        }

        \vspace{1.5cm}

        \rule{0.55\textwidth}{1pt}

        \vspace{0.7cm}

        {\large
        Competitive Programming
        }

        \vspace{0.3cm}

        {\large
        Data Structures \quad\textbullet\quad
        Algorithms \quad\textbullet\quad
        Mathematics
        }

        \vfill

        {\large
        gujiu6
        }

        % \vspace{0.3cm}

        % {\small
        % \today
        % }

        \vspace{1.8cm}

    \end{titlepage}
}

% ============================================================
% 正文开始
% ============================================================

\begin{document}

% ------------------------------------------------------------
% 封面
% ------------------------------------------------------------

\frontmatter

\MakeAlgorithmTitle

% ------------------------------------------------------------
% 前言
% ------------------------------------------------------------

\chapter*{前言}
\markboth{前言}{}

本手册以 \texttt{Kendieer/XCPC-Templates} 的 LaTeX 排版框架为基础，并将其中全部算法代码替换为 \texttt{gujiu6/Tem} 仓库的 C++ 模板。

\section*{使用说明}
\begin{itemize}[leftmargin=2em]
  \item 代码按基础算法、数据结构、图论、数学、线性代数与数值计算、字符串等专题编排；
  \item 所有代码均从 \texttt{code/} 目录直接读取，源码内容保持原仓库版本，不在 LaTeX 中二次改写；
  \item 算法的具体使用条件、复杂度和注释请以对应 C++ 源文件为准；
  \item 修改章节或代码文件后，请在项目根目录重新编译两遍，以刷新目录和页码。
\end{itemize}

\section*{编译方式}
需要支持中文的 XeLaTeX、\texttt{minted} 宏包以及 Python Pygments。于项目根目录执行：
\begin{verbatim}
xelatex -shell-escape -interaction=nonstopmode -halt-on-error main.tex
xelatex -shell-escape -interaction=nonstopmode -halt-on-error main.tex
\end{verbatim}

\begin{flushright}
    gujiu6\\
    \today
\end{flushright}


% ------------------------------------------------------------
% 目录（ctexbook 自动生成“目录”标题）
% ------------------------------------------------------------

\tableofcontents

\clearpage

% ------------------------------------------------------------
% 正文
% ------------------------------------------------------------

\mainmatter

% ============================================================
% 在这里引用章节（按需取消注释 / 新增 \include）
% ============================================================

\chapter{代码模板}

\section{竞赛主模板}

\inputminted{cpp}{code/template.cpp}

\chapter{基础算法}

\section{位运算}

\inputminted{cpp}{code/Basic/Bits.cpp}

\section{分治算法}

\inputminted{cpp}{code/Basic/DivideConquer.cpp}

\section{莫队算法}

\inputminted{cpp}{code/Basic/MoDui.cpp}

\section{前缀和与差分}

\inputminted{cpp}{code/Basic/PrefixAndDiff.cpp}

\section{常用工具}

\inputminted{cpp}{code/Basic/Tool.cpp}

\chapter{数据结构}

\section{平衡树与堆}

\subsection{堆}

\inputminted{cpp}{code/DateStructure/BBSTAndHeap/Heap.cpp}

\subsection{珂朵莉树（ODT）}

\inputminted{cpp}{code/DateStructure/BBSTAndHeap/ODT.cpp}

\subsection{Splay 平衡树}

\inputminted{cpp}{code/DateStructure/BBSTAndHeap/Splay.cpp}

\subsection{Treap 平衡树}

\inputminted{cpp}{code/DateStructure/BBSTAndHeap/Treap.cpp}

\section{树状数组}

\inputminted{cpp}{code/DateStructure/BIT/BIT.cpp}

\section{其他数据结构}

\subsection{并查集}

\inputminted{cpp}{code/DateStructure/Other/DSU.cpp}

\subsection{线性基}

\inputminted{cpp}{code/DateStructure/Other/LinearBasis.cpp}

\section{RMQ / ST 表}

\inputminted{cpp}{code/DateStructure/RMQ/RMQ.cpp}

\section{线段树}

\subsection{可持久化线段树}

\inputminted{cpp}{code/DateStructure/SegmentTree/PersistentSegmentTree.cpp}

\subsection{线段树}

\inputminted{cpp}{code/DateStructure/SegmentTree/SegmentTree1.cpp}

\section{Trie 字典树}

\inputminted{cpp}{code/DateStructure/TrieAndSet/Trie.cpp}


\chapter{图论}

\section{二分图匹配}

\inputminted{cpp}{code/GraphTheory/BipartiteGraph/Hungarian.cpp}

\section{强连通分量}

\inputminted{cpp}{code/GraphTheory/Connectivity/SCC.cpp}

\section{欧拉路径与欧拉回路}

\inputminted{cpp}{code/GraphTheory/EulerianGraph/Hierholzer.cpp}

\section{基础图论}

\subsection{图论基础}

\inputminted{cpp}{code/GraphTheory/Graph/Basic.cpp}

\subsection{最小生成树}

\inputminted{cpp}{code/GraphTheory/Graph/MinSpanningTree.cpp}

\subsection{最短路}

\inputminted{cpp}{code/GraphTheory/Graph/ShortestPath.cpp}

\section{最大流与最小割}

\inputminted{cpp}{code/GraphTheory/NetworkFlow/MaxFlowAndMinCut.cpp}

\section{树论}

\subsection{最近公共祖先（LCA）}

\inputminted{cpp}{code/GraphTheory/Tree/LCA.cpp}

\subsection{树的分治}

\inputminted{cpp}{code/GraphTheory/Tree/TreeCutting.cpp}

\subsection{树形动态规划}

\inputminted{cpp}{code/GraphTheory/Tree/TreeDP.cpp}

\subsection{树分治与 DSU on Tree}

\inputminted{cpp}{code/GraphTheory/Tree/TreeDivideandConquerAndDSUonTree.cpp}

\subsection{树论基础}

\inputminted{cpp}{code/GraphTheory/Tree/TreeTheory.cpp}


\chapter{线性代数与数值计算}

\section{线性方程组与数值算法}

\inputminted{cpp}{code/LinearAlgebraAndNumerical/LinearEquationsAndProgram.cpp}

\section{矩阵运算}

\inputminted{cpp}{code/LinearAlgebraAndNumerical/Matrix.cpp}

\chapter{数学}

\section{组合数学}

\subsection{组合数}

\inputminted{cpp}{code/Math/Comb/C(n,m).cpp}

\subsection{康托展开}

\inputminted{cpp}{code/Math/Comb/CanTor.cpp}

\section{数论}

\subsection{数论函数}

\inputminted{cpp}{code/Math/NumberTheory/ArithmeticFunctions.cpp}

\subsection{最大公约数与最小公倍数}

\inputminted{cpp}{code/Math/NumberTheory/GcdAndLcm.cpp}

\subsection{质数相关模板}

\inputminted{cpp}{code/Math/NumberTheory/Prime.cpp}

\section{多项式算法}

\inputminted{cpp}{code/Math/Polynomial/Polynomial.cpp}


\chapter{字符串算法}

\section{多模式串匹配}

\inputminted{cpp}{code/String/MultiplePatternMatch.cpp}

\section{字符串匹配}

\inputminted{cpp}{code/String/StringMatch.cpp}



% ============================================================
% 结束
% ============================================================

\backmatter

\end{document}